\documentclass[UTF8,a4paper,11pt]{ctexart}
\usepackage{amsmath,amssymb,bm,geometry,fancyhdr,enumitem,booktabs}
\geometry{left=22mm,right=22mm,top=22mm,bottom=22mm}
\pagestyle{fancy}
\fancyhf{}
\lhead{现代控制理论·第三周作业标准解答}
\rhead{习题七}
\cfoot{\thepage}
\setlength{\parindent}{2em}
\setlength{\parskip}{0.35em}
\begin{document}
\begin{center}
{\LARGE\bfseries 现代控制理论·第三周作业标准解答}\\[4pt]
{\large 习题七：7-13，7-19，7-20，7-21，7-27，7-30}
\end{center}

\section*{7-13}
已知
\[
\bm x(k+1)=A\bm x(k)+B u(k),\qquad
A=\begin{bmatrix}0&1\\-0.15&-0.8\end{bmatrix},\quad
B=\begin{bmatrix}1\\0\end{bmatrix},
\]
\[
\bm x(0)=\begin{bmatrix}0\\1\end{bmatrix},\qquad u(k)\equiv1.
\]
常值输入下先求平衡点 $\bm x_s$：
\[
\bm x_s=A\bm x_s+B,\qquad (I-A)\bm x_s=B.
\]
因此
\[
I-A=\begin{bmatrix}1&-1\\0.15&1.8\end{bmatrix},\qquad
\bm x_s=(I-A)^{-1}B
=\begin{bmatrix}12/13\\-1/13\end{bmatrix}.
\]
令 $\tilde{\bm x}(k)=\bm x(k)-\bm x_s$，则
\[
\tilde{\bm x}(k+1)=A\tilde{\bm x}(k),\qquad
\tilde{\bm x}(0)=
\begin{bmatrix}-12/13\\14/13\end{bmatrix}.
\]
矩阵 $A$ 的特征方程为
\[
\lambda^2+0.8\lambda+0.15=0,
\]
故
\[
\lambda_1=-0.3,\qquad \lambda_2=-0.5.
\]
相应特征向量可取
\[
\bm v_1=\begin{bmatrix}1\\-0.3\end{bmatrix},\qquad
\bm v_2=\begin{bmatrix}1\\-0.5\end{bmatrix}.
\]
由
\[
\tilde{\bm x}(0)=\frac{40}{13}\bm v_1-4\bm v_2
\]
可得
\[
\tilde{\bm x}(k)=\frac{40}{13}(-0.3)^k\bm v_1-4(-0.5)^k\bm v_2.
\]
所以
\[
\boxed{
\bm x(k)=
\begin{bmatrix}
\dfrac{12}{13}+\dfrac{40}{13}(-0.3)^k-4(-0.5)^k\\[8pt]
-\dfrac{1}{13}+\dfrac{40}{13}(-0.3)^{k+1}-4(-0.5)^{k+1}
\end{bmatrix}}
\]
其中 $k=0,1,2,\ldots$。

\section*{7-19}
连续部分
\[
G(s)=\frac{2}{s(0.1s+1)}=\frac{20}{s(s+10)}.
\]
其脉冲响应为
\[
g(t)=\mathcal L^{-1}\{G(s)\}=2(1-e^{-10t}).
\]
采样周期 $T=0.1\,\mathrm{s}$，令
\[
a=e^{-10T}=e^{-1}=0.367879.
\]
则开环脉冲传递函数
\[
\begin{aligned}
G(z)
&=\mathcal Z\{g(kT)\}\\
&=2\left(\frac{z}{z-1}-\frac{z}{z-a}\right)\\
&=\frac{2(1-a)z}{(z-1)(z-a)}.
\end{aligned}
\]
单位负反馈闭环脉冲传递函数为
\[
\Phi(z)=\frac{G(z)}{1+G(z)}
=\frac{2(1-a)z}{z^2+(1-3a)z+a}.
\]
代入 $a=e^{-1}$：
\[
\boxed{\Phi(z)=\frac{1.264241z}{z^2-0.103638z+0.367879}}
\]
单位阶跃输入
\[
X(z)=\frac{z}{z-1},
\]
所以
\[
Y(z)=\frac{2(1-a)z^2}{(z-1)\,[z^2+(1-3a)z+a]}.
\]
可整理为
\[
Y(z)=\frac{z}{z-1}
-\frac{z(z-a)}{z^2-(3a-1)z+a}.
\]
闭环特征根为
\[
z_{1,2}=0.051819\pm j0.604313=\rho e^{\pm j\theta},
\]
其中
\[
\rho=\sqrt a=e^{-1/2}=0.606531,\qquad
\theta=1.485257\ \mathrm{rad}.
\]
因此采样时刻的单位阶跃响应为
\[
\boxed{
y(kT)=1+\rho^k\left[-\cos(k\theta)+0.523008\sin(k\theta)\right],
\quad k=0,1,2,\ldots}
\]
前几个采样值为
\[
\begin{array}{c|cccccccc}
k&0&1&2&3&4&5&6&7\\ \hline
y(kT)&0&1.2642&1.3953&0.9438&0.8488&1.0050&1.0562&1.0040
\end{array}
\]
且 $y(kT)\to1$。

\section*{7-20}
连续部分为
\[
G(s)=\frac{K}{s(T_1s+1)},\qquad T_1>0,
\]
脉冲响应
\[
g(t)=K(1-e^{-t/T_1}).
\]
令
\[
a=e^{-T/T_1},\qquad 0<a<1,
\]
则
\[
G(z)=\frac{K(1-a)z}{(z-1)(z-a)}.
\]
闭环特征方程
\[
1+G(z)=0
\]
即
\[
P(z)=z^2+[K(1-a)-(1+a)]z+a=0.
\]
为使用劳斯判据，作双线性变换
\[
w=\frac{z-1}{z+1},\qquad z=\frac{1+w}{1-w}.
\]
单位圆内部对应 $w$ 平面的左半平面。代入 $P(z)$ 并乘去 $(1-w)^2$，得
\[
\begin{aligned}
P_w(w)=&\,[2(1+a)-K(1-a)]w^2\\
&+2(1-a)w+K(1-a).
\end{aligned}
\]
二阶多项式的劳斯表第一列为
\[
\begin{array}{c|c}
w^2&2(1+a)-K(1-a)\\
w^1&2(1-a)\\
w^0&K(1-a)
\end{array}
\]
因为 $0<a<1$，稳定要求第一列同号且取正，故
\[
K>0,\qquad 2(1+a)-K(1-a)>0.
\]
所以
\[
\boxed{
0<K<\frac{2(1+a)}{1-a}
=\frac{2(1+e^{-T/T_1})}{1-e^{-T/T_1}}}
\]
即题给条件正是闭环稳定条件。

\section*{7-21}
开环脉冲传递函数
\[
G(z)=\frac{K(z-0.5)}{(z-1)^2}.
\]
单位负反馈闭环特征方程
\[
1+G(z)=0
\]
即
\[
(z-1)^2+K(z-0.5)=0,
\]
整理为
\[
\boxed{z^2+(K-2)z+1-\frac K2=0}.
\]
根轨迹的稳定边界是单位圆。令单位圆上的交点
\[
z=e^{j\theta}.
\]
代入特征方程并乘以 $e^{-j\theta}$：
\[
e^{j\theta}+(K-2)+\left(1-\frac K2\right)e^{-j\theta}=0.
\]
取虚部：
\[
\frac K2\sin\theta=0.
\]
对 $K>0$，有
\[
\sin\theta=0,
\]
所以根轨迹只能在 $z=1$ 或 $z=-1$ 穿越单位圆。

在 $z=1$ 时，代入特征方程得到 $K=0$，这是根轨迹起点；在 $z=-1$ 时
\[
1-(K-2)+1-\frac K2=0,
\]
故
\[
4-\frac{3K}{2}=0.
\]
因此正增益下的临界值为
\[
\boxed{K_{\mathrm{cr}}=\frac83}.
\]
在该点闭环极点为
\[
z=-1,\qquad z=\frac13.
\]
所以闭环稳定范围为
\[
\boxed{0<K<\frac83}.
\]

\section*{7-27}
闭环传递函数
\[
\Phi_B(z)=\frac{0.368z+0.264}{z^2-z+0.632}.
\]
若采用单位负反馈，则误差与输入之间满足
\[
\frac{E(z)}{R(z)}=1-\Phi_B(z).
\]
因此
\[
\begin{aligned}
1-\Phi_B(z)
&=\frac{z^2-z+0.632-(0.368z+0.264)}{z^2-z+0.632}\\
&=\frac{z^2-1.368z+0.368}{z^2-z+0.632}\\
&=\frac{(z-1)(z-0.368)}{z^2-z+0.632}.
\end{aligned}
\]
闭环极点为 $0.5\pm j0.6181$，模小于 1，可使用离散终值定理。

\subsection*{(1) $r(t)=1(t)$}
采样后
\[
R(z)=\frac{z}{z-1}.
\]
所以
\[
\begin{aligned}
e_{ss}
&=\lim_{z\to1}\frac{z-1}{z}E(z)\\
&=\lim_{z\to1}[1-\Phi_B(z)]=0.
\end{aligned}
\]
\[
\boxed{e_{ss}=0}
\]

\subsection*{(2) $r(t)=t$}
采样周期记为 $T$，有
\[
R(z)=\frac{Tz}{(z-1)^2}.
\]
于是
\[
\begin{aligned}
e_{ss}
&=\lim_{z\to1}\frac{z-1}{z}
\frac{(z-1)(z-0.368)}{z^2-z+0.632}
\frac{Tz}{(z-1)^2}\\
&=T\frac{1-0.368}{1-1+0.632}=T.
\end{aligned}
\]
\[
\boxed{e_{ss}=T}
\]
若采用归一化采样周期 $T=1$，则 $e_{ss}=1$。

\subsection*{(3) $r(t)=t^2/2$}
\[
R(z)=\frac{T^2z(z+1)}{2(z-1)^3}.
\]
误差传递函数只含一个因子 $(z-1)$，代入终值定理后仍剩一个 $(z-1)$ 在分母，因此
\[
\boxed{e_{ss}=\infty}.
\]

综上：
\[
\boxed{1(t):\ 0,\qquad t:\ T,\qquad t^2/2:\ \infty}.
\]

\section*{7-30}
采样周期 $T=1$，连续部分
\[
G_0(s)=\frac{1}{s(s+1)}.
\]
其脉冲响应为
\[
g_0(t)=1-e^{-t}.
\]
令
\[
a=e^{-1},
\]
则被控对象脉冲传递函数
\[
\begin{aligned}
G(z)
&=\mathcal Z\{1-a^k\}\\
&=\frac{z}{z-1}-\frac{z}{z-a}\\
&=\frac{(1-a)z}{(z-1)(z-a)}.
\end{aligned}
\]
写成 $z^{-1}$ 形式：
\[
G(z)=\frac{(1-a)z^{-1}}
{(1-z^{-1})(1-az^{-1})}.
\]
对象本身至少有一拍延迟，因此单位阶跃的最少拍闭环响应应取
\[
\boxed{\Phi(z)=z^{-1}}.
\]
此时采样输出为
\[
c(0)=0,\qquad c(k)=1\quad(k\ge1),
\]
即一拍后精确到达给定值，且稳态误差为零。

由
\[
\Phi(z)=\frac{D(z)G(z)}{1+D(z)G(z)}
\]
可得
\[
D(z)=\frac{\Phi(z)}{G(z)[1-\Phi(z)]}.
\]
代入 $\Phi(z)=z^{-1}$ 与 $G(z)$：
\[
\begin{aligned}
D(z)
&=\frac{z^{-1}}
{\dfrac{(1-a)z^{-1}}{(1-z^{-1})(1-az^{-1})}(1-z^{-1})}\\
&=\frac{1-az^{-1}}{1-a}.
\end{aligned}
\]
因此最少拍数字控制器为
\[
\boxed{
D(z)=\frac{1-e^{-1}z^{-1}}{1-e^{-1}}
=\frac{z-e^{-1}}{(1-e^{-1})z}}
\]
数值形式为
\[
\boxed{D(z)=1.58198-0.58198z^{-1}}.
\]
校验：
\[
D(z)G(z)=\frac{z^{-1}}{1-z^{-1}},
\qquad
\frac{D(z)G(z)}{1+D(z)G(z)}=z^{-1},
\]
故系统在一拍内结束过渡过程，且单位阶跃稳态误差为零。

\end{document}
